\section{A uniform obstruction to repairing the 193-state runtime by register allocation}
\label{sec:runtime193-all-layouts}

The supplied runtime core identifies native control \(301\) with \(88\),
redirects every target \(301\) to \(88\), removes the two rows of \(301\),
and installs
\[
 (88,0)\mapsto(0,L,51),\qquad (51,0)\mapsto(0,R,122).
\]
Its 193-state count does not make it a blank-start consistency machine.

Let
\[
 \mathrm{SK}=(51,36,25,48,32,33,22,55,60).
\]
For any permutation \(\pi\) of the nine physical register positions, first
apply the register allocation to the primary 194-state table: every
entry edge from a non-arithmetic-box control to \(\mathrm{SK}_j\) is
redirected to \(\mathrm{SK}_{\pi(j)}\), while the arithmetic-box rows are
retained. Then apply the runtime fold above. The order matters:
the new helper edge \(88\mathbin{/}0\to51\) is installed after allocation
and remains a literal edge to \(51\).

\paragraph{Proposition.}
Every one of these \(9!=362{,}880\) physical layouts, started in native
state \(74\) on the blank tape, reaches an infinite rightward writing
ray. In particular, none of these reallocations repairs the startup
of this runtime core.

\paragraph{Proof.}
The first 78 transitions do not use any of the routing rows changed by
the physical allocation. The complete common finite prefix ends at
\[
 t=78,\quad q=45,\quad h=4,\quad
             \{i:\mathrm{tape}(i)=1\}=\{-3,1,4\}.
\]
The full trace and the check that no allocation-dependent row is used
are included in the accompanying certificate.

Put \(r=\pi(7)\). The next transition is
\((45,1)\mapsto(0,R,\mathrm{SK}_{r})\).
It removes the one at cell \(4\), enters \(\mathrm{SK}_{r}\) at head
position \(5\), and leaves the entire suffix from position \(5\) blank.
For \(1\le j\le8\), the unchanged skip chain satisfies
\[
  (\mathrm{SK}_{j},0)\mapsto(0,R,\mathrm{SK}_{j-1}).
\]
After precisely \(r\) further transitions, the configuration therefore
has state \(51\), head \(5+r\), and one-cells exactly \(\{-3,1\}\).
This is time \(79+r\).

The following six zero-reading transitions are fixed in every layout:
\[
\begin{array}{c|c}
\text{current state} & \text{read-zero action}\\ \hline
51  & 0R122\\
122 & 0L137\\
137 & 0R30\\
30  & 0R45\\
45  & 0R24\\
24  & 0L35
\end{array}
\]
They reach state \(35\), head \(7+r\), at time \(85+r\), still with
one-cells exactly \(\{-3,1\}\). Finally,
\[
                 (35,0)\mapsto(1,R,35).
\]
The head continually enters a fresh blank cell, writes one, and moves
right again. Thus it never leaves state \(35\) and never halts.
All arguments are finite case checks followed by elementary induction,
so the nonhalting conclusion for each fixed table is provable in PA.
\hfill\(\square\)

The executable search also explicitly enumerates the \(8!=40{,}320\)
layouts fixing physical position \(6\). Every one reaches the
hypothesis of the six-step ray lemma, with no timeout or unclassified
case. Their literal prefixes have a total of \(3{,}336{,}480\) transitions.
The time to state \(51\) is \(79+\pi(7)\), giving 5,040 cases each at
times \(79,80,81,82,83,84,86,87\). Time \(85\) is absent in that subset
because \(\pi(7)\ne6\).

This is an obstruction for the specified runtime fold, allocation
transformation, and native start. It does not rule out a different
bootstrap, another runtime representation, another 193-state design,
or a still smaller PA-consistency machine. In particular, it is not a
state-count lower bound.
